% =====================================================================
%  Seção de Avaliação do NEXUS Trust Score — gerada a partir de
%  scripts/avaliar_score.py sobre 631 empresas reais (231 sancionadas +
%  400 controles ativos-limpos). Ground truth: sanções oficiais vigentes
%  (CEIS/CNEP/Lista Suja do Trabalho Escravo).
%
%  Pacotes necessários no preâmbulo:
%     \usepackage{booktabs}
%     \usepackage{threeparttable}   % para as notas de rodapé da tabela
% =====================================================================

\section{Avaliação da Qualidade da Detecção}
\label{sec:avaliacao}

Para além do desempenho computacional, avaliamos a \emph{qualidade da
detecção} do NEXUS Trust Score como um problema de classificação binária.
A pontuação de risco de cada empresa é definida como $\text{risco}=100-\text{score}$
e usada como a saída contínua do classificador. O \textit{ground truth} é
composto por empresas com sanção oficial vigente (rótulo positivo,
CEIS/CNEP/Lista Suja do Trabalho Escravo) e por uma amostra de controle de
empresas ativas sem sanção e sem dívida ativa na União (rótulo negativo).
A amostra avaliada contém $n=631$ empresas, sendo $231$ positivas
($36{,}6\%$ de prevalência). Todos os intervalos entre colchetes são
intervalos de confiança de $95\%$ obtidos por \textit{bootstrap}
($2{,}000$ reamostragens).

\begin{table}[t]
  \centering
  \begin{threeparttable}
  \caption{Discriminação do NEXUS Trust Score frente a baselines
           (limiar de decisão: risco $\geq 50$).}
  \label{tab:avaliacao}
  \begin{tabular}{lccc}
    \toprule
    Modelo & ROC-AUC & PR-AUC & F1 \\
    \midrule
    NEXUS (completo)\tnote{a}     & $1{,}000$ & $1{,}000$ & $1{,}000$ \\
    NEXUS (apenas inferido)\tnote{b} & $0{,}396$ & $0{,}508$ & $0{,}362$ \\
    \midrule
    Baseline: dívida ativa        & $0{,}755$ & ---       & $0{,}676$ \\
    Baseline: situação cadastral  & $0{,}610$ & ---       & $0{,}362$ \\
    Baseline: aleatório           & $0{,}490$ & $0{,}349$ & $0{,}431$ \\
    \bottomrule
  \end{tabular}
  \begin{tablenotes}\footnotesize
    \item[a] Inclui as camadas confirmadas por fonte oficial. Como o rótulo
      (sanção) é também um insumo do score, há \emph{vazamento de rótulo}:
      este valor é um teto otimista, não uma medida de poder preditivo.
    \item[b] Camadas oficiais (sanção e dívida) removidas; restam apenas os
      sinais inferidos (cadastro da Receita Federal, contágio da rede
      societária e cruzamento de endereços).
    \item[] PR-AUC dos baselines omitido: como os controles foram amostrados
      sem dívida nem irregularidade cadastral, esses sinais valem zero em
      todos os negativos por construção, tornando a métrica degenerada.
  \end{tablenotes}
  \end{threeparttable}
\end{table}

A Tabela~\ref{tab:avaliacao} revela um resultado que o benchmark de
desempenho não captura. No modo \emph{completo}, o NEXUS reproduz
perfeitamente o \textit{ground truth} (ROC-AUC $=1{,}000$), mas esse valor
reflete \textbf{vazamento de rótulo}: a sanção que define a classe positiva
é um dos insumos diretos do score. Trata-se de uma verificação de
sanidade --- o motor de fato ancora seu veredito nos fatos oficiais --- e
não de evidência de capacidade preditiva.

A avaliação cientificamente relevante é a do \emph{motor inferido}, com as
camadas oficiais desligadas. Nesse cenário, a discriminação cai para
ROC-AUC $=0{,}396$ (IC $95\%$: $[0{,}342;\,0{,}448]$), \emph{abaixo do
acaso}, com recall de apenas $0{,}221$: os sinais de rede societária e
co-localização, isoladamente, não separam empresas sancionadas das
empresas-controle. Isso \textbf{não} implica que tais sinais sejam inúteis,
mas sim que o \textit{ground truth} de sanções é inadequado para validá-los:
sanções administrativas medem um construto de risco (idoneidade em
licitações, corrupção, trabalho escravo) distinto daquele que a análise
relacional busca --- estruturas de laranjas e empresas de fachada. Por fim,
dois baselines triviais (presença de dívida ativa e situação cadastral
negativa) já alcançam ROC-AUC de $0{,}755$ e $0{,}610$, indicando que boa
parte do poder de separação observável vem de fatos cadastrais simples.

\paragraph{Limitações e trabalhos futuros.} Estes resultados expõem a
principal lacuna de validação: a ausência de um conjunto rotulado de
\emph{fraude relacional confirmada} (esquemas de laranjas e fachada
judicialmente caracterizados). A validação adequada do motor inferido --- o
diferencial do NEXUS --- exige essa base rotulada, idealmente revisada
manualmente, sobre a qual as métricas da Tabela~\ref{tab:avaliacao} devem ser
recomputadas. O arcabouço de avaliação (\texttt{scripts/avaliar\_score.py})
fica disponível para essa reprodução.
